% Записи notation-glossary (SSOT). Код: hL, hT → $a$, $\Delta t$; planck cell → $V_P$ ($v$ — скорость).
% Описания — в \newcommand (не inline в description={...}: xkeyval и `}`).
% Имена макросов: нельзя \NotationDescc0 (TeX читает как \NotationDescc+0); только \NotationDescCzero.
% glossaries раскрывает description в преамбуле: математика только \(...\) (нижние индексы
% внутри math: \(p_0\), не ASCII p\_0 в прозе). Не $...$ и не \Cref.
% Колонка «Обозначение» (поле name): только символы и индексы в имени ($N_{\mathrm{ring}}$,
% $N_{12}$), без подстановки численных значений (${=}512$, ${=}13$ и т.п.).

\newcommand{\NotationDescMlayer}{решётка, локальный закон $g$.}
\newcommand{\NotationDescTlayer}{осреднённая физика, readout, согласование с экспериментом.}
\newcommand{\NotationDescLambda}{счётное множество узлов носителя.}
\newcommand{\NotationDescVp}{планковская ячейка: объём $V_P=\ell_P^3$ на M (не скорость $v$). См. \hyperref[sec:carrier-object]{§0}, \hyperref[sec:planck-lattice]{§1}.}
\newcommand{\NotationDescg}{локальный обратимый закон эволюции; конфигурация в такт t. См. \hyperref[sec:descent-ca]{§0}.}
\newcommand{\NotationDesca}{элементарный шаг решётки (шаг--такт); в нат.\ ед.\ на $M$: $a{=}1$; код \texttt{hL}. См. \hyperref[sec:planck-lattice]{§1}.}
\newcommand{\NotationDescdt}{длительность одного такта $g$ (шаг--такт); в нат.\ ед.\ на $M$: $\Delta t{=}1$; код \texttt{hT}. См. \hyperref[sec:coords-dh-dt]{§0.1}.}
\newcommand{\NotationDesctau}{номер такта (поколение CA); связь с \(t_{\mathrm{SI}}\) после SI-моста.}
\newcommand{\NotationDescCzero}{тактовая скорость сигнала; \(c_0=a/\Delta t\). См. \hyperref[sec:two-c]{§1}.}
\newcommand{\NotationDescCspeed}{скорость света после SI-моста: $c=\kappa c_0$. См. \hyperref[sec:two-c]{§1}.}
\newcommand{\NotationDesckgeom}{безразмерный коэффициент геометрии носителя (kappa в коде). См. \hyperref[sec:fcc-kappa]{§1}.}
\newcommand{\NotationDescdhdt}{внешние координаты dh и dt (dh=a, dt=\(\Delta t\)). См. \hyperref[sec:coords-dh-dt]{§0.1}.}
\newcommand{\NotationDescszero}{ступень действия на одну ячейку и один такт; порождает лестницу квантов. См. \hyperref[sec:theorem-51]{теорему~5.1}.}
\newcommand{\NotationDescpzero}{импульсный квант: \(p_0=s_0/dh\). См. \hyperref[sec:theorem-51]{§5.2.1}.}
\newcommand{\NotationDescEzero}{энергетический квант: \(E_0=s_0/\Delta t\). См. \hyperref[sec:theorem-51]{§5.2.1}.}
\newcommand{\NotationDescFzero}{силовой квант: \(F_0=s_0/(dh\cdot\Delta t)\). См. \hyperref[sec:theorem-51]{§5.2.1}.}
\newcommand{\NotationDescmzero}{массовый квант лестницы: \(m_0=E_0/c_0^2\); SI-мост. См. \hyperref[sec:theorem-51]{§5.2.1}, \hyperref[sec:full-ladder]{§6}.}
\newcommand{\NotationDescL}{макроскопическая координата длины (readout).}
\newcommand{\NotationDesclambda}{характерная длина на шаге шкалы (комптоновская, боровская, \ldots).}
\newcommand{\NotationDesclambdaC}{комптоновская длина электрона.}
\newcommand{\NotationDesclambdak}{шаги IR-лестницы.}
\newcommand{\NotationDescellP}{планковская длина; на носителе \(\ell_P=a\) (SI-мост). См. \hyperref[sec:planck-lattice]{§1}.}
\newcommand{\NotationDescNc}{безразмерное \(\lambda_C/\ell_P\); не «голое» N.}
\newcommand{\NotationDescloglambda}{порядок длины в таблице шкалы.}
\newcommand{\NotationDesctSI}{время прибора (осреднение по тактам).}
\newcommand{\NotationDesctP}{планковское время; \(\Delta t=\kappa\,t_P\) (SI-мост).}
\newcommand{\NotationDesclogtau}{порядок числа тактов (логарифмическая шкала).}
\newcommand{\NotationDescNring}{размер фазового кольца; \(N_{\mathrm{ring}}=2^{\lfloor B_V\rfloor}\) (дискретная фаза). См. \hyperref[thm:cell-bit-budget]{теорему о бит-бюджете}.}
\newcommand{\NotationDescBhv}{информационный бюджет ячейки $B_V$; связь с \(N_{\mathrm{ring}}\) (дискретная фаза). См. \hyperref[sec:planckon-internal-spectrum]{§5.0.4}.}
\newcommand{\NotationDescphiPhi}{дискретная фаза $\phaseVarphi_{\mathrm{disc}}$ и импульс за тик $\symCapitalPhi$.}
\newcommand{\NotationDescNphi}{число секторов Heisenberg-кванта фазы.}
\newcommand{\NotationDescNFCC}{число ближайших соседей FCC; степени \(N_{12}^m\) --- геометрия звезды.}
\newcommand{\NotationDescNpack}{число упаковки атомной ячейки (открытый вывод).}
\newcommand{\NotationDescNhier}{число ступеней ослабления на T (VEV).}

\newcommand{\NotationSecCarrierTitle}{\textit{Слои и носитель}}
\newcommand{\NotationSecLatticeTitle}{\textit{Решётка: длина и время (не путать с \(N_{\mathrm{ring}}\))}}
\newcommand{\NotationSecLadderTitle}{\textit{Лестница квантов на \((dh,dt)\)}}
\newcommand{\NotationSecLengthTitle}{\textit{Длина по шкале}}
\newcommand{\NotationSecTimeTitle}{\textit{Время}}
\newcommand{\NotationSecCarrierTTitle}{\textit{Макроуровень}}
\newcommand{\NotationSecPhaseTitle}{\textit{Фаза, регистр ячейки (логарифмы в основе 2)}}
\newcommand{\NotationSecNamedNTitle}{\textit{Именованные N (не ось длины)}}

\newnotationsectionM{symb:sec-carrier}{carrier-00}{\NotationSecCarrierTitle}
\newnotationentryM{symb:Mlayer}{carrier-01a}{\ensuremath{M}}{\NotationDescMlayer}{\NotationUnitDash}
\newnotationentryM{symb:Lambda}{carrier-02}{\ensuremath{\Lambda}}{\NotationDescLambda}{\NotationUnitDash}
\newnotationentryM{symb:Vp}{carrier-03}{\ensuremath{V_P}}{\NotationDescVp}{\NotationUnitMetreCubed}
\newnotationentryM{symb:g}{carrier-04}{\ensuremath{g{,}\,\Psi^t}}{\NotationDescg}{\NotationUnitDash}

\newnotationsectionT{symb:sec-carrier-t}{carrier-t-00}{\NotationSecCarrierTTitle}
\newnotationentryT{symb:Tlayer}{carrier-01b}{\ensuremath{T}}{\NotationDescTlayer}{\NotationUnitDash}

\newnotationsectionM{symb:sec-lattice}{lattice-00}{\NotationSecLatticeTitle}
\newnotationentryM{symb:a}{lattice-01}{\ensuremath{a}}{\NotationDesca}{\NotationUnitMetreMcode}
\newnotationentryM{symb:dt}{lattice-02}{\ensuremath{\Delta t}}{\NotationDescdt}{\NotationUnitSecondDtcode}
\newnotationentryM{symb:tau}{lattice-03}{\ensuremath{\tau\in\mathbb{Z}}}{\NotationDesctau}{\NotationUnitTick}
\newnotationentryM{symb:c0}{lattice-04}{\ensuremath{c_0}}{\NotationDescCzero}{\NotationUnitMetrePerSecond}
\newnotationentryM{symb:kgeom}{lattice-06}{\ensuremath{\kappa}}{\NotationDesckgeom}{\NotationUnitDash}
\newnotationentryM{symb:dhdt}{lattice-07}{\ensuremath{(dh{,}dt)}}{\NotationDescdhdt}{\NotationUnitDash}
\newnotationentryT{symb:c}{lattice-05}{\ensuremath{c}}{\NotationDescCspeed}{\NotationUnitMetrePerSecond}

\newnotationsectionM{symb:sec-ladder}{ladder-00}{\NotationSecLadderTitle}
\newnotationentryM{symb:s0}{ladder-01}{\ensuremath{s_0}}{\NotationDescszero}{\NotationUnitJouleSecond}
\newnotationentryM{symb:p0}{ladder-02}{\ensuremath{p_0}}{\NotationDescpzero}{\NotationUnitKgMetrePerSecond}
\newnotationentryM{symb:E0}{ladder-03}{\ensuremath{E_0}}{\NotationDescEzero}{\NotationUnitJoule}
\newnotationentryM{symb:F0}{ladder-04}{\ensuremath{F_0}}{\NotationDescFzero}{\NotationUnitNewton}
\newnotationentryM{symb:m0}{ladder-05}{\ensuremath{m_0}}{\NotationDescmzero}{\NotationUnitKilogram}

\newnotationsectionT{symb:sec-length}{length-00}{\NotationSecLengthTitle}
\newnotationentryT{symb:L}{length-01}{\ensuremath{L}}{\NotationDescL}{\NotationUnitMetre}
\newnotationentryT{symb:lambda}{length-02}{\ensuremath{\lambda}}{\NotationDesclambda}{\NotationUnitMetre}
\newnotationentryT{symb:lambdaC}{length-03}{\ensuremath{\lambda_C}}{\NotationDesclambdaC}{\NotationUnitMetre}
\newnotationentryT{symb:lambdak}{length-04}{\ensuremath{\lambda_k}}{\NotationDesclambdak}{\NotationUnitMetre}
\newnotationentryT{symb:ellP}{length-05}{\ensuremath{\ell_P}}{\NotationDescellP}{\NotationUnitMetre}
\newnotationentryT{symb:Nc}{length-06}{\ensuremath{N_c}}{\NotationDescNc}{\NotationUnitDash}
\newnotationentryT{symb:loglambda}{length-07}{\ensuremath{\log_{10}(\lambda/\ell_P)}}{\NotationDescloglambda}{\NotationUnitDash}

\newnotationsectionT{symb:sec-time}{time-00}{\NotationSecTimeTitle}
\newnotationentryT{symb:tSI}{time-01}{\ensuremath{t_{\mathrm{SI}}}}{\NotationDesctSI}{\NotationUnitSecond}
\newnotationentryT{symb:tP}{time-02}{\ensuremath{t_P}}{\NotationDesctP}{\NotationUnitSecond}
\newnotationentryT{symb:logtau}{time-03}{\ensuremath{\log_{10}|\tau|}}{\NotationDesclogtau}{\NotationUnitDash}

\newnotationsectionM{symb:sec-phase}{phase-00}{\NotationSecPhaseTitle}
\newnotationentryM{symb:Nring}{phase-01}{\ensuremath{N_{\mathrm{ring}}}}{\NotationDescNring}{\NotationUnitDash}
\newnotationentryM{symb:BV}{phase-02}{\ensuremath{B_V}}{\NotationDescBhv}{\NotationUnitDash}
\newnotationentryM{symb:phiPhi}{phase-03}{\ensuremath{\varphi_{\mathrm{disc}}{,}\,\Phi}}{\NotationDescphiPhi}{\NotationUnitPhaseTick}
\newnotationentryM{symb:Nphi}{phase-04}{\ensuremath{N_\phi}}{\NotationDescNphi}{\NotationUnitDash}

\newnotationsectionM{symb:sec-namedN}{namedN-00}{\NotationSecNamedNTitle}
\newnotationentryM{symb:N12}{namedN-01}{\ensuremath{N_{12}}}{\NotationDescNFCC}{\NotationUnitDash}
\newnotationentryM{symb:Npack}{namedN-02}{\ensuremath{N_{\mathrm{pack}}}}{\NotationDescNpack}{\NotationUnitDash}
\newnotationentryT{symb:Nhier}{namedN-03}{\ensuremath{N_{\mathrm{hier}}}}{\NotationDescNhier}{\NotationUnitDash}

% Макро-физика (Максвелл, Ньютон, КМ, …) — тултипы в формулах через макросы.

\newcommand{\NotationSecOperatorsTitle}{\textit{Операторы}}
\newcommand{\NotationSecMacroPhysicsTitle}{\textit{Макроскопическая физика (поля, константы, SI)}}
\newcommand{\NotationDescNabla}{векторный дифференциальный оператор
$(\partial/\partial x,\,\partial/\partial y,\,\partial/\partial z)$;
на скаляр $\varphi$ даёт \emph{градиент} $\nabla\varphi$.}
\newcommand{\NotationDescDiv}{дивергенция векторного поля $\mathbf{A}$; $\nabla\cdot\mathbf{A}$
(уравнение непрерывности, Гаусса).}
\newcommand{\NotationDescCurl}{ротор векторного поля $\mathbf{A}$; $\nabla\times\mathbf{A}$
(уравнения Фарадея и Ампера--Максвелла).}
\newcommand{\NotationDescLaplacian}{оператор Лапласа $\nabla^2=\nabla\cdot\nabla$
(уравнение Пуассона, член $-\frac{\hbar^2}{2m}\nabla^2$ в КМ).}
\newcommand{\NotationDescPartialT}{производная по времени $\partial/\partial t$ (локально на $T$);
размерность задаётся операндом.}
\newcommand{\NotationDescEfld}{вектор электрического поля.}
\newcommand{\NotationDescBfld}{вектор магнитной индукции.}
\newcommand{\NotationDescJfld}{плотность электрического тока.}
\newcommand{\NotationDescRhoQ}{объёмная плотность заряда (уравнение Гаусса).}
\newcommand{\NotationDescRhoM}{массовая плотность (гидродинамика, не заряд).}
\newcommand{\NotationDescEpsZero}{электрическая постоянная вакуума \(\varepsilon_0\).}
\newcommand{\NotationDescMuZero}{магнитная постоянная вакуума \(\mu_0\).}
\newcommand{\NotationDescVfld}{вектор скорости.}
\newcommand{\NotationDescRfld}{радиус-вектор точки.}
\newcommand{\NotationDescFfld}{сила (механика Ньютона).}
\newcommand{\NotationDescPfld}{давление (гидродинамика).}
\newcommand{\NotationDescSentropy}{энтропия.}
\newcommand{\NotationDescTtemp}{температура.}
\newcommand{\NotationDescHham}{гамильтониан / оператор Гамильтона.}
\newcommand{\NotationDescPsi}{волновая функция.}
\newcommand{\NotationDescHbar}{постоянная Планка \(\hbar\).}
\newcommand{\NotationDescEcharge}{элементарный заряд \(e\).}
\newcommand{\NotationDescMe}{масса электрона \(m_e\).}
\newcommand{\NotationDescMmass}{масса тела (Ньютон).}
\newcommand{\NotationDescVarphi}{фаза (Маделунг, \(\psi=\sqrt{\rho}\,e^{i\varphi/\hbar}\)).}
\newcommand{\NotationDescClight}{скорость света в уравнениях поля и механики.}
\newcommand{\NotationDesckB}{постоянная Больцмана \(k_B\).}
\newcommand{\NotationDescGgrav}{гравитационная постоянная \(G\).}
\newcommand{\NotationDescVecP}{вектор импульса \(\mathbf{p}\) (пространственная часть \(p^\mu\)).}

\newnotationsectionT{symb:sec-operators}{macro-00}{\NotationSecOperatorsTitle}
\newnotationentryT{symb:nabla}{macro-01}{\ensuremath{\nabla}}{\NotationDescNabla}{\NotationUnitDash}
\newnotationentryT{symb:div}{macro-02}{\ensuremath{\nabla\cdot}}{\NotationDescDiv}{\NotationUnitDash}
\newnotationentryT{symb:curl}{macro-03}{\ensuremath{\nabla\times}}{\NotationDescCurl}{\NotationUnitDash}
\newnotationentryT{symb:laplacian}{macro-04}{\ensuremath{\nabla^2}}{\NotationDescLaplacian}{\NotationUnitDash}
\newnotationentryT{symb:partialt}{macro-05}{\ensuremath{\partial_t}}{\NotationDescPartialT}{\NotationUnitDash}

\newnotationsectionT{symb:sec-macrophys}{macrophys-00}{\NotationSecMacroPhysicsTitle}
\newnotationentryT{symb:Efld}{macrophys-01}{\ensuremath{\mathbf{E}}}{\NotationDescEfld}{\NotationUnitVoltPerMetre}
\newnotationentryT{symb:Bfld}{macrophys-02}{\ensuremath{\mathbf{B}}}{\NotationDescBfld}{\NotationUnitTesla}
\newnotationentryT{symb:Jfld}{macrophys-03}{\ensuremath{\mathbf{j}}}{\NotationDescJfld}{\NotationUnitAmperePerMetreSq}
\newnotationentryT{symb:rhoQ}{macrophys-04}{\ensuremath{\rho}}{\NotationDescRhoQ}{\NotationUnitCoulombPerMetreCubed}
\newnotationentryT{symb:rhoM}{macrophys-05}{\ensuremath{\rho}}{\NotationDescRhoM}{\NotationUnitKgPerMetreCubed}
\newnotationentryT{symb:eps0}{macrophys-06}{\ensuremath{\varepsilon_0}}{\NotationDescEpsZero}{\NotationUnitFaradPerMetre}
\newnotationentryT{symb:mu0}{macrophys-07}{\ensuremath{\mu_0}}{\NotationDescMuZero}{\NotationUnitHenryPerMetre}
\newnotationentryT{symb:Vfld}{macrophys-08}{\ensuremath{\mathbf{v}}}{\NotationDescVfld}{\NotationUnitMetrePerSecond}
\newnotationentryT{symb:Rfld}{macrophys-09}{\ensuremath{\mathbf{r}}}{\NotationDescRfld}{\NotationUnitMetre}
\newnotationentryT{symb:Ffld}{macrophys-10}{\ensuremath{\mathbf{F}}}{\NotationDescFfld}{\NotationUnitNewton}
\newnotationentryT{symb:Pfld}{macrophys-11}{\ensuremath{p}}{\NotationDescPfld}{\NotationUnitPascal}
\newnotationentryT{symb:Sentropy}{macrophys-12}{\ensuremath{S}}{\NotationDescSentropy}{\NotationUnitJoulePerKelvin}
\newnotationentryT{symb:Ttemp}{macrophys-13}{\ensuremath{T}}{\NotationDescTtemp}{\NotationUnitKelvin}
\newnotationentryT{symb:Hham}{macrophys-14}{\ensuremath{H}}{\NotationDescHham}{\NotationUnitJoule}
\newnotationentryT{symb:Psi}{macrophys-15}{\ensuremath{\psi}}{\NotationDescPsi}{\NotationUnitWavefunction}
\newnotationentryT{symb:hbar}{macrophys-16}{\ensuremath{\hbar}}{\NotationDescHbar}{\NotationUnitJouleSecond}
\newnotationentryT{symb:echar}{macrophys-17}{\ensuremath{e}}{\NotationDescEcharge}{\NotationUnitCoulomb}
\newnotationentryT{symb:me}{macrophys-18}{\ensuremath{m_e}}{\NotationDescMe}{\NotationUnitKilogram}
\newnotationentryT{symb:Mmass}{macrophys-19}{\ensuremath{m}}{\NotationDescMmass}{\NotationUnitKilogram}
\newnotationentryT{symb:varphi}{macrophys-20}{\ensuremath{\varphi}}{\NotationDescVarphi}{\NotationUnitJoule}
\newnotationentryT{symb:clight}{macrophys-21}{\ensuremath{c}}{\NotationDescClight}{\NotationUnitMetrePerSecond}
\newnotationentryT{symb:kB}{macrophys-22}{\ensuremath{k_B}}{\NotationDesckB}{\NotationUnitJoulePerKelvin}
\newnotationentryT{symb:Ggrav}{macrophys-23}{\ensuremath{G}}{\NotationDescGgrav}{\NotationUnitGravConstant}
\newnotationentryT{symb:vecP}{macrophys-24}{\ensuremath{\mathbf{p}}}{\NotationDescVecP}{\NotationUnitKgMetrePerSecond}

% Autogen-глоссарий символов (TeX SSOT; правки руками, без mass-wire по главам).
\newcommand{\NotationSecAutogenTitle}{\textit{Символы в формулах (autogen)}}
\newnotationsection{symb:sec-autogen}{autogen-00}{\NotationSecAutogenTitle}
\newcommand{\NotationDescAutovecX}{вектор $\mathbf{x}$ (координата / узел).}
\newglossaryentry{symb:auto:vecX}{%
  type=notation, sort={autogen-vecX},
  name={\ensuremath{\mathbf{x}}},
  description={\NotationDescAutovecX}%
}
\newcommand{\NotationDescAutovecY}{вектор $\mathbf{y}$.}
\newglossaryentry{symb:auto:vecY}{%
  type=notation, sort={autogen-vecY},
  name={\ensuremath{\mathbf{y}}},
  description={\NotationDescAutovecY}%
}
\newcommand{\NotationDescAutovecZ}{вектор $\mathbf{z}$ (спинор / регистр).}
\newglossaryentry{symb:auto:vecZ}{%
  type=notation, sort={autogen-vecZ},
  name={\ensuremath{\mathbf{z}}},
  description={\NotationDescAutovecZ}%
}
\newcommand{\NotationDescAutopartialK}{производная по координате $k$ ($\partial_k$).}
\newglossaryentry{symb:auto:partialK}{%
  type=notation, sort={autogen-partialK},
  name={\ensuremath{\partial_k}},
  description={\NotationDescAutopartialK}%
}
\newcommand{\NotationDescAutoCapitalPsi}{$\Psi$ (капитальная; поле на решётке).}
\newglossaryentry{symb:auto:CapitalPsi}{%
  type=notation, sort={autogen-CapitalPsi},
  name={\ensuremath{\Psi}},
  description={\NotationDescAutoCapitalPsi}%
}
\newcommand{\NotationDescAutoalphaFs}{константа тонкой структуры $\alpha_{\mathrm{fs}}$.}
\newglossaryentry{symb:auto:alphaFs}{%
  type=notation, sort={autogen-alphaFs},
  name={\ensuremath{\alpha_{\mathrm{fs}}}},
  description={\NotationDescAutoalphaFs}%
}
\newcommand{\NotationDescAutostressSigmaIk}{компонента тензора напряжений $\sigma_{ik}$.}
\newglossaryentry{symb:auto:stressSigmaIk}{%
  type=notation, sort={autogen-stressSigmaIk},
  name={\ensuremath{\sigma_{ik}}},
  description={\NotationDescAutostressSigmaIk}%
}
\newcommand{\NotationDescAutoRthree}{трёхмерное евклидово пространство $\mathbb{R}^3$.}
\newglossaryentry{symb:auto:Rthree}{%
  type=notation, sort={autogen-Rthree},
  name={\ensuremath{\mathbb{R}^3}},
  description={\NotationDescAutoRthree}%
}
\newcommand{\NotationDescAutoRfour}{пространство-время $\mathbb{R}^4$ (континуальный предел).}
\newglossaryentry{symb:auto:Rfour}{%
  type=notation, sort={autogen-Rfour},
  name={\ensuremath{\mathbb{R}^4}},
  description={\NotationDescAutoRfour}%
}
\newcommand{\NotationDescAutoR}{множество действительных чисел $\mathbb{R}$.}
\newglossaryentry{symb:auto:R}{%
  type=notation, sort={autogen-R},
  name={\ensuremath{\mathbb{R}}},
  description={\NotationDescAutoR}%
}
\newcommand{\NotationDescAutoC}{множество комплексных чисел $\mathbb{C}$.}
\newglossaryentry{symb:auto:C}{%
  type=notation, sort={autogen-C},
  name={\ensuremath{\mathbb{C}}},
  description={\NotationDescAutoC}%
}
\newcommand{\NotationDescAutoZ}{целые числа $\mathbb{Z}$ (такт, фаза).}
\newglossaryentry{symb:auto:Z}{%
  type=notation, sort={autogen-Z},
  name={\ensuremath{\mathbb{Z}}},
  description={\NotationDescAutoZ}%
}
\newcommand{\NotationDescAutocalB}{оператор осреднения $\mathcal{B}$ (макроуровень).}
\newglossaryentry{symb:auto:calB}{%
  type=notation, sort={autogen-calB},
  name={\ensuremath{\mathcal{B}}},
  description={\NotationDescAutocalB}%
}
\newcommand{\NotationDescAutocalC}{пространство конфигураций $\mathcal{C}$.}
\newglossaryentry{symb:auto:calC}{%
  type=notation, sort={autogen-calC},
  name={\ensuremath{\mathcal{C}}},
  description={\NotationDescAutocalC}%
}
\newcommand{\NotationDescAutocalN}{нормировка / число состояний $\mathcal{N}$.}
\newglossaryentry{symb:auto:calN}{%
  type=notation, sort={autogen-calN},
  name={\ensuremath{\mathcal{N}}},
  description={\NotationDescAutocalN}%
}
\newcommand{\NotationDescAutovarepsilon}{символ $\varepsilon$ (диэлектрическая проницаемость в контексте).}
\newglossaryentry{symb:auto:varepsilon}{%
  type=notation, sort={autogen-varepsilon},
  name={\ensuremath{\varepsilon}},
  description={\NotationDescAutovarepsilon}%
}
\newcommand{\NotationDescAutophi}{угол / фаза $\phi$.}
\newglossaryentry{symb:auto:phi}{%
  type=notation, sort={autogen-phi},
  name={\ensuremath{\phi}},
  description={\NotationDescAutophi}%
}
\newcommand{\NotationDescAutoPhi}{фазовый поток / капитальная $\Phi$.}
\newglossaryentry{symb:auto:Phi}{%
  type=notation, sort={autogen-Phi},
  name={\ensuremath{\Phi}},
  description={\NotationDescAutoPhi}%
}
\newcommand{\NotationDescAutoalpha}{безразмерный параметр $\alpha$.}
\newglossaryentry{symb:auto:alpha}{%
  type=notation, sort={autogen-alpha},
  name={\ensuremath{\alpha}},
  description={\NotationDescAutoalpha}%
}
\newcommand{\NotationDescAutobeta}{параметр $\beta$.}
\newglossaryentry{symb:auto:beta}{%
  type=notation, sort={autogen-beta},
  name={\ensuremath{\beta}},
  description={\NotationDescAutobeta}%
}
\newcommand{\NotationDescAutogamma}{параметр $\gamma$ (не путать с $\Gamma$).}
\newglossaryentry{symb:auto:gamma}{%
  type=notation, sort={autogen-gamma},
  name={\ensuremath{\gamma}},
  description={\NotationDescAutogamma}%
}
\newcommand{\NotationDescAutoGamma}{$\Gamma$ (группа / капитальная гамма).}
\newglossaryentry{symb:auto:Gamma}{%
  type=notation, sort={autogen-Gamma},
  name={\ensuremath{\Gamma}},
  description={\NotationDescAutoGamma}%
}
\newcommand{\NotationDescAutodelta}{вариация / $\delta$-функция в контексте.}
\newglossaryentry{symb:auto:delta}{%
  type=notation, sort={autogen-delta},
  name={\ensuremath{\delta}},
  description={\NotationDescAutodelta}%
}
\newcommand{\NotationDescAutoDelta}{оператор Лапласа / приращение $\Delta$.}
\newglossaryentry{symb:auto:Delta}{%
  type=notation, sort={autogen-Delta},
  name={\ensuremath{\Delta}},
  description={\NotationDescAutoDelta}%
}
\newcommand{\NotationDescAutokappa}{безразмерный $\kappa$ (геометрия; не $\caKappaGeom$ без контекста).}
\newglossaryentry{symb:auto:kappa}{%
  type=notation, sort={autogen-kappa},
  name={\ensuremath{\kappa}},
  description={\NotationDescAutokappa}%
}
\newcommand{\NotationDescAutomu}{параметр $\mu$ (химический потенциал / магнитный момент в контексте).}
\newglossaryentry{symb:auto:mu}{%
  type=notation, sort={autogen-mu},
  name={\ensuremath{\mu}},
  description={\NotationDescAutomu}%
}
\newcommand{\NotationDescAutonu}{частота / параметр $\nu$.}
\newglossaryentry{symb:auto:nu}{%
  type=notation, sort={autogen-nu},
  name={\ensuremath{\nu}},
  description={\NotationDescAutonu}%
}
\newcommand{\NotationDescAutopi}{число $\pi$.}
\newglossaryentry{symb:auto:pi}{%
  type=notation, sort={autogen-pi},
  name={\ensuremath{\pi}},
  description={\NotationDescAutopi}%
}
\newcommand{\NotationDescAutorho}{скаляр $\rho$ (плотность: масса, заряд или $|\psi|^2$ — по контексту).}
\newglossaryentry{symb:auto:rho}{%
  type=notation, sort={autogen-rho},
  name={\ensuremath{\rho}},
  description={\NotationDescAutorho}%
}
\newcommand{\NotationDescAutosigma}{$\sigma$ (проводимость, напряжение, сечение — по контексту).}
\newglossaryentry{symb:auto:sigma}{%
  type=notation, sort={autogen-sigma},
  name={\ensuremath{\sigma}},
  description={\NotationDescAutosigma}%
}
\newcommand{\NotationDescAutotau}{время / номер такта $\tau$.}
\newglossaryentry{symb:auto:tau}{%
  type=notation, sort={autogen-tau},
  name={\ensuremath{\tau}},
  description={\NotationDescAutotau}%
}
\newcommand{\NotationDescAutotheta}{угол $\theta$.}
\newglossaryentry{symb:auto:theta}{%
  type=notation, sort={autogen-theta},
  name={\ensuremath{\theta}},
  description={\NotationDescAutotheta}%
}
\newcommand{\NotationDescAutoomega}{частота $\omega$.}
\newglossaryentry{symb:auto:omega}{%
  type=notation, sort={autogen-omega},
  name={\ensuremath{\omega}},
  description={\NotationDescAutoomega}%
}
\newcommand{\NotationDescAutoOmega}{$\Omega$ (объём / капитальная омега).}
\newglossaryentry{symb:auto:Omega}{%
  type=notation, sort={autogen-Omega},
  name={\ensuremath{\Omega}},
  description={\NotationDescAutoOmega}%
}
\newcommand{\NotationDescAutoell}{длина $\ell$ (не $\ell_P$ без индекса).}
\newglossaryentry{symb:auto:ell}{%
  type=notation, sort={autogen-ell},
  name={\ensuremath{\ell}},
  description={\NotationDescAutoell}%
}
\newcommand{\NotationDescAutozeta}{параметр $\zeta$.}
\newglossaryentry{symb:auto:zeta}{%
  type=notation, sort={autogen-zeta},
  name={\ensuremath{\zeta}},
  description={\NotationDescAutozeta}%
}
\newcommand{\NotationDescAutoenergyE}{энергия $E$ (скаляр, не поле $\mathbf{E}$).}
\newglossaryentry{symb:auto:energyE}{%
  type=notation, sort={autogen-energyE},
  name={\ensuremath{E}},
  description={\NotationDescAutoenergyE}%
}


\notationBlkMcloseblock
\notationBlkTcloseblock
